\chapter*{Preface}
\label{chap:preface}
\addcontentsline{toc}{chapter}{Preface}

The first three volumes of the project are technical gauges. Volume One
reads the project as a mathematical theorem. Volume Two reads it as a
physical instrument. Volume Three reads it as a compiler. Each of those
volumes asks the reader to enter through a discipline: mathematical
patience, physical intuition, algorithmic imagination. Volume Four, the
reference trace, asks for trust in a transcript. Each of those volumes
assumes a reader willing to learn a register before reading the argument
in it.

This volume asks for ordinary language.

The argument the three technical volumes carry is the same. A record can
be made. The record can be distinguished, admitted, counted, compared,
represented, observed, held open, staged, executed, scaled, witnessed,
localized, transported, halted, measured, compiled, and inferred. Each
verb is the name of a public obligation. In the Lean code that supports
the project, each obligation is also the name of a class. The classes
are not implementation labels. They are verdicts. Each name marks a
condition that the argument had to satisfy before the next move was
admissible. A move that fails its named obligation does not pass the
gate the obligation guards. The argument proceeds only along paths that
the named obligations have permitted.

The class names in the code, taken in order, form a public argument. A
reader who does not know any mathematics, physics, or compiler theory
can follow the procession of names. The names exert their own pressure.
\texttt{DISTINGUISHABLE} is not decoration. It says that if you cannot
tell it apart, it cannot enter the ledger. \texttt{ADMISSIBLE} is not
decoration. It says that if it cannot enter under rule, it is not yet a
reading. \texttt{REPEATABLE} is not decoration. It says that if it
cannot happen again under the same conditions, it is not yet an
instrument. The names refuse weak claims by their syntax. They make
evasion harder. They make the obligations inspectable. A reader who
has not learned the formal machinery can still notice that a claim has
dropped one of these names along the way, and can still ask the speaker
to restore the missing obligation before the argument continues.

The class names of the project expose their obligations almost
brutally. \texttt{GUNGAN} is a rude name for the function of holding
awkward states without resolving them. \texttt{BULLSHIT} is a rude name
for the function of measuring overhead that has become attached to a
record. \texttt{ACOLYTE} is a rude name for the carrier who repeats a
claim because it belongs to a system rather than because the record has
earned it. The names are not polished. The polish would let the
obligation hide. \texttt{ACCUMULATED\_OPERATIONAL\_OVERHEAD\_INDEX}
could mean the same thing as \texttt{BULLSHIT}. The first form would
let the cost slip behind a clinical phrase. The second form would not.
The project kept the second form.

Two example registers carry the argument. The first is the pre-flight
checklist that precedes a commercial takeoff. The checklist is brief,
named, and disciplined, and it lives inside an apparatus the reader
already trusts. A passenger who has never read a mathematical paper
still trusts that the captain's checklist controls the takeoff. The
trust is not arbitrary. It is the result of decades of accident
analysis that drove the airlines to adopt named checklists in place of
cockpit improvisation. The trust the passenger gives the checklist is
the same trust the reader is asked to give the procession of class
names: not faith in the names themselves, but recognition that the
named discipline has held under pressure long enough to be relied on.

The second register is a used-car inspection before purchase. The
inspection moves through the full procession of class names from first
mark to final inference. A buyer can use the inspection's structure to
say, publicly and without cheating, whether the car is worth buying at
the offered price. The decision is local. It does not involve
regulators, public-health authorities, or institutional politics. It
involves a buyer, a seller, a mechanic, and a sum of money that the
buyer can choose to spend or keep. The decision is the right size for
the rhetorical gauge to work on. The procession is large enough to
require all 37 class names. The decision is small enough that no class
name is asked to do work it cannot earn.

The procession ends at \texttt{INFERRED}, not at \texttt{TRUTH},
because the discipline of the procession is to license a public sentence
under a record, not to claim possession of a fact.

The class names include some that are deliberately impolite.
\texttt{GUNGAN}, \texttt{BULLSHIT}, \texttt{PROPAGANDA},
\texttt{ACOLYTE}, \texttt{FINITE\_ELEPHANT}. The names were not adopted
to be clever. They were adopted because polished alternatives let the
obligations hide. A reader who finds the names jarring may ask whether
the jarring is the name's failure or the name's success. A name that
should jar but does not has been polished out of its work. A name that
jars and reveals nothing has been adopted for its shock alone. The
project's satirical names jar and reveal. They identify a specific real
function inside the record that no polished name carries as well.
\texttt{BULLSHIT} is a name for overhead that has attached to a record
without contributing to it. The name is honest.

The class names are quoted in code font when they are named, but their
formal definitions live in the other volumes. A reader who wants to see
what \texttt{DISTINGUISHABLE} looks like as a typeclass should consult
Volume Three.

The first question is small:

\begin{center}
\emph{What kind of claim can a name carry?}
\end{center}

Everything that follows is the consequence of refusing to answer that
question carelessly.
